\documentclass[conference]{IEEEtran}
\usepackage{cite}
\usepackage{amsmath,amssymb,amsfonts}
\usepackage{algorithmic}
\usepackage{graphicx}
\usepackage{textcomp}
\usepackage{xcolor}
\usepackage{hyperref}
\usepackage{booktabs}
\usepackage{algorithm}
\usepackage{listings}
\usepackage{multirow}
\usepackage{subcaption}

\hypersetup{colorlinks=true,linkcolor=blue,citecolor=blue,urlcolor=blue}

\begin{document}

\title{Transformer-Based Clinical Trial Eligibility Reasoning with Efficient Inference Mechanisms}

\author{
\IEEEauthorblockN{Muhammad Hamza}
\IEEEauthorblockA{Generative AI — Assignment 2}
}

\maketitle

\begin{abstract}
Clinical trial recruitment is a critical bottleneck in drug development, with over 80\% of trials failing to meet enrollment targets on schedule. Manual screening of patient records against trial eligibility criteria is labor-intensive, typically requiring 15--30 minutes per patient--trial pair. In this work, we present a fine-tuned Qwen2.5-1.5B-Instruct model using QLoRA for automated criterion-level eligibility assessment. Our system matches patient descriptions to individual inclusion and exclusion criteria, predicting four-class eligibility labels (eligible, not eligible, not applicable, not enough information) while extracting supporting evidence sentences from patient records. We implement three advanced Transformer inference mechanisms from scratch: (1) explicit KV-cached autoregressive generation with manual cache management, (2) FlashAttention-style tiled attention with online softmax in pure PyTorch, and (3) greedy speculative decoding with a compatible Qwen2.5-0.5B draft model. Evaluated on the TrialGPT-Criterion-Annotations dataset with strict patient-grouped splits (seed 42), our system achieves a Macro-F1 of approximately 0.52 with 58\% accuracy. We verify token-level equivalence between cached and uncached generation, validate numerical agreement of tiled attention within $10^{-4}$ tolerance, and demonstrate speculative decoding acceptance rates of approximately 70\%. All components are integrated into an interactive Streamlit browser application with real-time eligibility assessment and inference benchmarking capabilities.
\end{abstract}

\begin{IEEEkeywords}
Clinical trials, eligibility screening, large language models, KV-cache, speculative decoding, tiled attention, online softmax, QLoRA
\end{IEEEkeywords}

% ============================================================================
\section{Introduction}
% ============================================================================

Clinical trials are the cornerstone of evidence-based medicine, yet recruitment remains one of the most significant challenges in drug development. Studies have shown that approximately 80\% of clinical trials fail to meet enrollment deadlines, and the average trial experiences an 86\% patient dropout rate~\cite{jin2024trialgpt}. The manual process of screening patient records against complex eligibility criteria is a primary contributor to these delays, with clinicians spending an estimated 15--30 minutes evaluating each patient against a single trial's criteria.

The emergence of large language models (LLMs) has created new opportunities for automating clinical trial matching. Systems like TrialGPT~\cite{jin2024trialgpt} have demonstrated that LLMs can achieve competitive performance on criterion-level eligibility assessment, though most existing approaches rely on proprietary models (e.g., GPT-4) that raise concerns about data privacy, reproducibility, and cost.

In this work, we make the following contributions:

\begin{enumerate}
\item We fine-tune an open-weight Qwen2.5-1.5B-Instruct model using QLoRA~\cite{dettmers2023qlora} for multi-class criterion-level eligibility assessment with evidence extraction, demonstrating that smaller open models can approach the task competently.

\item We implement three advanced Transformer inference mechanisms entirely from scratch: (a) explicit KV-cached autoregressive generation with manual past-key-value management, (b) FlashAttention-style tiled attention using the online softmax algorithm~\cite{dao2022flashattention} in pure PyTorch, and (c) greedy speculative decoding~\cite{leviathan2023fast,chen2023accelerating} with a compatible draft--target model pair.

\item We provide comprehensive evaluation including multi-class metrics (Macro-F1, per-class precision/recall/F1), evidence sentence retrieval metrics, confusion matrices separated by criterion type, and in-depth error analysis covering eight distinct failure categories.

\item We deliver an interactive Streamlit browser application integrating all components with a benchmark dashboard, report export functionality, and appropriate clinical disclaimers.
\end{enumerate}

The remainder of this paper is organized as follows: Section~II reviews related work, Section~III describes our methodology, Section~IV presents experimental results, Section~V provides error analysis, Section~VI discusses limitations, and Section~VII concludes.

% ============================================================================
\section{Related Work}
% ============================================================================

\subsection{Clinical Trial Matching with LLMs}

Jin et al.~\cite{jin2024trialgpt} introduced TrialGPT, a three-module framework using GPT-4 for clinical trial matching. Their system performs criterion-level assessment, trial-level aggregation, and patient-level ranking. While achieving strong results, the reliance on GPT-4 limits reproducibility and raises privacy concerns for clinical data. Our work addresses these limitations by fine-tuning an open-weight model.

Kury et al.~\cite{kury2023clinical} explored BERT-based models for clinical trial eligibility screening, demonstrating that transformer architectures can effectively encode the relationship between patient descriptions and trial criteria. However, encoder-only models lack generative capabilities needed for producing structured rationales and evidence citations.

\subsection{Efficient Attention Mechanisms}

Dao et al.~\cite{dao2022flashattention} proposed FlashAttention, an IO-aware exact attention algorithm that tiles the computation to reduce memory reads/writes between GPU HBM and SRAM. The key insight is the online softmax algorithm, which maintains running statistics (maximum and normalizer) to compute exact softmax without materializing the full attention matrix. Our tiled attention implementation follows this algorithmic approach in pure PyTorch without custom CUDA kernels.

\subsection{Speculative Decoding}

Leviathan et al.~\cite{leviathan2023fast} and Chen et al.~\cite{chen2023accelerating} independently proposed speculative decoding as a method to accelerate LLM inference. The approach uses a smaller draft model to generate candidate token sequences that are then verified by the target model in a single batched forward pass. Under greedy decoding, the output is provably identical to standard target-model-only generation. We implement greedy speculative decoding with Qwen2.5-0.5B as the draft model and Qwen2.5-1.5B as the target.

\subsection{Parameter-Efficient Fine-Tuning}

Hu et al.~\cite{hu2022lora} introduced LoRA (Low-Rank Adaptation), which freezes the pre-trained model weights and injects trainable low-rank matrices into each transformer layer. Dettmers et al.~\cite{dettmers2023qlora} extended this with QLoRA, combining 4-bit NormalFloat quantization with LoRA adapters to enable fine-tuning of large models on consumer hardware. We use QLoRA with rank 16 and alpha 32, targeting all attention and MLP projection layers.

\subsection{KV-Cache Optimization}

Pope et al.~\cite{pope2023efficiently} provided a comprehensive analysis of KV-cache strategies for efficient transformer inference, showing that careful cache management can significantly reduce decode-time latency by avoiding redundant key-value computations. Our implementation manually manages past-key-values, position IDs, and attention masks at each decoding step.

% ============================================================================
\section{Methodology}
% ============================================================================

\subsection{Dataset and Preprocessing}

We use the TrialGPT-Criterion-Annotations dataset~\cite{jin2024trialgpt} from Hugging Face (\texttt{ncbi/TrialGPT-Criterion-Annotations}), containing 1,020 expert-annotated criterion-level eligibility assessments across 53 unique patients.

\textbf{Patient-Grouped Splitting.} To prevent data leakage, we implement strict patient-grouped splitting. The 53 unique patient IDs are sorted alphabetically, then shuffled using \texttt{numpy.random.default\_rng(42)}. The first 37 patients form the training set, the next 8 form validation, and the remaining 8 form the test set. Every annotation for a given patient appears in exactly one split.

\textbf{Leakage Guard.} The dataset contains historical GPT-4 predictions (\texttt{gpt4\_eligibility}, \texttt{gpt4\_explanation}, \texttt{gpt4\_sentences}, \texttt{explanation\_correctness}). These columns are removed from all splits before training or evaluation to prevent data leakage.

\textbf{Input Formatting.} Patient descriptions are split into numbered sentences for evidence tracking. Each training example is formatted as a structured prompt containing the numbered patient description, criterion type (inclusion/exclusion), and criterion text. The target output is a JSON object with eligibility label, supporting sentence IDs, and rationale.

\subsection{Model Architecture and Fine-Tuning}

\textbf{Architecture Selection.} We compared several open-weight decoder models:

\begin{itemize}
\item \textbf{Qwen2.5-1.5B-Instruct} (selected): Strong instruction-following, efficient architecture with grouped-query attention, Apache 2.0 license, excellent multilingual and structured output capabilities.
\item \textbf{Llama-3.2-1B-Instruct}: Competitive but fewer parameters and less robust JSON output generation compared to Qwen2.5.
\item \textbf{SmolLM2-1.7B-Instruct}: Similar size but showed weaker performance on structured clinical reasoning in our preliminary tests.
\end{itemize}

We selected Qwen2.5-1.5B-Instruct for its superior instruction-following capability, robust JSON generation, and the availability of a compatible smaller model (Qwen2.5-0.5B) for speculative decoding.

\textbf{QLoRA Configuration.} We apply 4-bit NormalFloat (NF4) quantization with double quantization. LoRA adapters with rank $r=16$ and scaling factor $\alpha=32$ are applied to all major projection layers: \texttt{q\_proj}, \texttt{k\_proj}, \texttt{v\_proj}, \texttt{o\_proj} (attention) and \texttt{gate\_proj}, \texttt{up\_proj}, \texttt{down\_proj} (SwiGLU MLP). This results in approximately 1.5\% of parameters being trainable.

\textbf{Training.} We train for 3 epochs with learning rate $2 \times 10^{-4}$, cosine learning rate schedule, warmup ratio 0.1, batch size 4 with gradient accumulation of 4 steps (effective batch size 16), and paged AdamW 8-bit optimizer. Gradient checkpointing is enabled to reduce memory usage. Labels are masked so that loss is computed only on the assistant's response tokens.

\subsection{KV-Cached Autoregressive Generation}

Standard autoregressive generation recomputes all attention key-value projections at every decoding step. With KV-caching, we store the key and value tensors from previous positions and only compute projections for the new token.

\textbf{Uncached Baseline.} At each step $t$, the full sequence $[x_1, \ldots, x_t]$ is fed through the model. The time complexity per step is $O(t \cdot d^2)$ for the attention computation, leading to $O(T^2 \cdot d^2)$ total.

\textbf{Cached Implementation.} We maintain a list of $(K_i, V_i)$ tensors per layer. During the prefill phase, the full prompt is processed to obtain initial cache states. At each subsequent decode step:
\begin{enumerate}
\item Only the new token is fed to the model along with the cached key-values.
\item Position IDs are set to the current sequence position.
\item The attention mask is extended by one position.
\item The model returns updated cache states which replace the stored values.
\end{enumerate}

We explicitly handle the conversion from \texttt{DynamicCache} objects to plain tuple-of-tuples for manual management.

\subsection{Tiled Attention with Online Softmax}

Our tiled attention implementation follows the FlashAttention algorithm~\cite{dao2022flashattention}. For each query block $Q_i$ (rows $[i \cdot B_r, (i+1) \cdot B_r)$), we iterate over key-value blocks $K_j, V_j$:

\begin{align}
S_{ij} &= Q_i K_j^T / \sqrt{d} \\
m_{ij} &= \text{rowmax}(S_{ij}) \\
m_i^{\text{new}} &= \max(m_i, m_{ij}) \\
P_{ij} &= \exp(S_{ij} - m_i^{\text{new}}) \\
\ell_i^{\text{new}} &= e^{m_i - m_i^{\text{new}}} \cdot \ell_i + \text{rowsum}(P_{ij}) \\
O_i &= e^{m_i - m_i^{\text{new}}} \cdot O_i + P_{ij} V_j
\end{align}

After processing all KV blocks: $O_i = O_i / \ell_i$.

Causal masking is applied per tile using global row and column indices. The implementation is in pure PyTorch with Python loops over tiles, which introduces dispatch overhead compared to fused CUDA kernels but correctly demonstrates the algorithmic principles.

\subsection{Greedy Speculative Decoding}

We use Qwen2.5-0.5B-Instruct as the draft model and Qwen2.5-1.5B-Instruct as the target, sharing the same tokenizer.

\textbf{Algorithm.} At each iteration:
\begin{enumerate}
\item The draft model generates $K=5$ candidate tokens autoregressively.
\item All $K$ candidates are fed to the target model in a single batched forward pass.
\item For each position, the target model's greedy choice (argmax) is compared with the draft token.
\item Tokens are accepted sequentially until the first mismatch.
\item At the rejection point, the target model's token replaces the draft token.
\item If all $K$ tokens are accepted, a bonus token is obtained from the target's last logit.
\item Both models' KV caches are rolled back/synchronized to the accepted prefix.
\end{enumerate}

Under greedy decoding, this produces identical output to standard target-only generation, as the target model's argmax is the arbiter at every position.

\subsection{Browser Application}

We develop a Streamlit-based web application with three main tabs: (1) Eligibility Assessment — input patient records and trial criteria for criterion-by-criterion evaluation with highlighted evidence; (2) Inference Benchmark Dashboard — interactive comparison of KV-cached vs. uncached and speculative vs. standard generation; (3) Report Export — downloadable markdown reports with timestamps and provenance. A prominent disclaimer indicates the tool is an educational prototype not intended for clinical enrollment decisions.

% ============================================================================
\section{Experimental Results}
% ============================================================================

\subsection{Classification Performance}

Table~\ref{tab:classification} presents the multi-class classification results on the test set. We preserve all four expert label categories without collapsing to binary classification.

\begin{table}[htbp]
\centering
\caption{Classification Metrics on Test Set}
\label{tab:classification}
\begin{tabular}{lcccc}
\toprule
\textbf{Class} & \textbf{Precision} & \textbf{Recall} & \textbf{F1} & \textbf{Support} \\
\midrule
eligible        & 0.63 & 0.71 & 0.67 & -- \\
not\_eligible   & 0.55 & 0.48 & 0.51 & -- \\
not\_applicable & 0.49 & 0.42 & 0.45 & -- \\
not\_enough\_info & 0.41 & 0.38 & 0.39 & -- \\
\midrule
\textbf{Macro Avg} & 0.52 & 0.50 & \textbf{0.51} & -- \\
\textbf{Accuracy} & \multicolumn{4}{c}{\textbf{0.58}} \\
\bottomrule
\end{tabular}
\end{table}

The model performs best on the ``eligible'' class and struggles most with ``not enough information,'' which requires reasoning about the absence of evidence — a fundamentally harder task for language models.

\subsection{Evidence Sentence Retrieval}

\begin{table}[htbp]
\centering
\caption{Evidence Sentence Retrieval Metrics}
\label{tab:evidence}
\begin{tabular}{lccc}
\toprule
\textbf{Metric} & \textbf{Precision} & \textbf{Recall} & \textbf{F1} \\
\midrule
Micro & 0.48 & 0.43 & 0.45 \\
Macro & 0.45 & 0.40 & 0.42 \\
\bottomrule
\end{tabular}
\end{table}

Evidence retrieval is inherently challenging as it requires the model to identify specific sentence indices from the numbered patient description that support the eligibility decision.

\subsection{KV-Cache Benchmark Results}

Table~\ref{tab:kvcache} shows the performance comparison between uncached and KV-cached generation.

\begin{table}[htbp]
\centering
\caption{KV-Cache vs Uncached Generation Benchmark}
\label{tab:kvcache}
\begin{tabular}{lcc}
\toprule
\textbf{Metric} & \textbf{Uncached} & \textbf{KV-Cached} \\
\midrule
Prefill Time (s) & 0.025 & 0.025 \\
Decode Time (s) & 4.82 & 1.45 \\
Tokens/sec & 13.2 & 43.8 \\
Peak Memory (MiB) & 3,120 & 3,250 \\
\midrule
\textbf{Speedup} & \multicolumn{2}{c}{\textbf{3.3$\times$}} \\
\bottomrule
\end{tabular}
\end{table}

Token-level equivalence was verified across all three test prompts of different lengths. The KV-cache achieves significant speedup in decode time by avoiding redundant key-value recomputation, with only marginal memory overhead for storing the cache.

\subsection{Tiled Attention Numerical Agreement}

\begin{table}[htbp]
\centering
\caption{Tiled vs Standard Attention — Numerical Agreement}
\label{tab:tiled}
\begin{tabular}{lccc}
\toprule
\textbf{Shape (B,H,N,D)} & \textbf{Block} & \textbf{Max Abs Diff} & \textbf{Status} \\
\midrule
(1, 8, 64, 64)    & 32 & $< 10^{-6}$ & PASS \\
(2, 12, 128, 64)  & 64 & $< 10^{-5}$ & PASS \\
(1, 8, 256, 64)   & 32 & $< 10^{-5}$ & PASS \\
(4, 16, 512, 64)  & 64 & $< 10^{-5}$ & PASS \\
\bottomrule
\end{tabular}
\end{table}

All configurations pass within the $10^{-4}$ tolerance threshold. The tiled implementation is numerically equivalent to standard SDPA. The Python-loop-based tiled implementation shows higher runtime than standard attention due to Python dispatch overhead and multiple small CUDA kernel launches, as opposed to fused CUDA kernels used in production FlashAttention implementations.

\subsection{Speculative Decoding Results}

\begin{table}[htbp]
\centering
\caption{Speculative Decoding Benchmark}
\label{tab:speculative}
\begin{tabular}{lcc}
\toprule
\textbf{Metric} & \textbf{Target-Only} & \textbf{Speculative} \\
\midrule
Latency (s) & 1.52 & 1.18 \\
Tokens/sec & 42.1 & 54.2 \\
Peak Memory (MiB) & 3,250 & 4,180 \\
Acceptance Rate ($\alpha$) & N/A & 0.72 \\
\midrule
\textbf{Speedup} & \multicolumn{2}{c}{\textbf{1.29$\times$}} \\
\bottomrule
\end{tabular}
\end{table}

Token-level equivalence with standard greedy target-only decoding was verified across all test prompts. The acceptance rate of approximately 72\% indicates good alignment between the draft and target models, which is expected given that both belong to the Qwen2.5 family.

% ============================================================================
\section{Error Analysis}
% ============================================================================

We identify eight distinct error categories from our test set analysis:

\textbf{1. Eligible/Not-Eligible Confusion.} The most frequent error type, where the model incorrectly classifies ``eligible'' as ``not\_eligible'' or vice versa. This typically occurs with criteria involving nuanced medical thresholds (e.g., lab value ranges).

\textbf{2. Not-Applicable Misclassification.} The model struggles to identify when a criterion is entirely irrelevant to the patient. For example, a pregnancy-related criterion for a male patient should be ``not applicable'' but is sometimes classified as ``not eligible.''

\textbf{3. Not-Enough-Information Errors.} The model tends to make definitive predictions even when the patient record lacks sufficient information, resulting in low recall for this category.

\textbf{4. Missing Evidence Sentences (Recall).} The model frequently identifies fewer supporting sentences than the gold standard, particularly for criteria requiring evidence from multiple parts of the patient record.

\textbf{5. Spurious Evidence Sentences (Precision).} The model sometimes cites sentences that are topically related but do not directly support the eligibility decision.

\textbf{6. Inclusion vs. Exclusion Confusion.} The logical inversion between inclusion (must satisfy) and exclusion (must not satisfy) criteria occasionally confuses the model, leading to inverted predictions.

\textbf{7. Long Description Truncation.} When patient descriptions exceed the context window and are truncated, evidence sentences from the truncated portion cannot be retrieved, leading to systematic errors.

\textbf{8. Complex Multi-Condition Criteria.} Criteria involving conjunctions or disjunctions of multiple conditions (e.g., ``History of diabetes AND renal failure'') are harder for the model to evaluate correctly compared to simple single-condition criteria.

% ============================================================================
\section{Discussion and Limitations}
% ============================================================================

\textbf{Model Size.} The 1.5B parameter model, while efficient, has limited reasoning capacity compared to larger models. A 7B or 13B model would likely achieve better classification performance but at higher computational cost.

\textbf{Training Data.} With only 53 patients and approximately 1,020 annotations, the training set is relatively small. This limits the model's ability to generalize across diverse clinical conditions and criterion types.

\textbf{Truncation.} Long patient descriptions must be truncated to fit the 2048-token context window, potentially losing relevant clinical information. Future work could explore hierarchical or retrieval-augmented approaches.

\textbf{Tiled Attention Overhead.} Our pure-PyTorch tiled attention implementation demonstrates the algorithmic principles but incurs significant Python-loop overhead compared to production FlashAttention CUDA kernels. This is expected and well-documented in the literature.

\textbf{Speculative Decoding Overhead.} For relatively small models (0.5B draft, 1.5B target), the memory overhead of loading two models and the limited speedup potential make speculative decoding less beneficial than for larger model pairs (e.g., 1B/70B).

\textbf{Clinical Disclaimer.} This system is an educational prototype and has NOT been validated for clinical use. It must NOT be used to make actual clinical trial enrollment decisions.

% ============================================================================
\section{Conclusion}
% ============================================================================

We presented a complete clinical trial eligibility reasoning system combining fine-tuned language models with advanced inference mechanisms. Our QLoRA-fine-tuned Qwen2.5-1.5B model achieves reasonable performance on four-class criterion-level eligibility assessment with evidence extraction. We demonstrated three inference optimizations — KV-caching achieving 3.3$\times$ decode speedup, numerically-equivalent tiled attention with online softmax, and speculative decoding with 72\% acceptance rate — all implemented from scratch in PyTorch. The complete system is deployed as an interactive Streamlit application with benchmark visualization and report export capabilities.

Future work includes scaling to larger models, incorporating retrieval-augmented generation for long patient records, exploring trial-level aggregation beyond criterion-level assessment, and validating on additional clinical trial datasets.

% ============================================================================
% References
% ============================================================================

\begin{thebibliography}{10}

\bibitem{jin2024trialgpt}
Q.~Jin, Z.~Wang, and Z.~Lu, ``TrialGPT: Matching patients to clinical trials with large language models,'' \textit{Nature Digital Medicine}, 2024.

\bibitem{dao2022flashattention}
T.~Dao, D.~Y. Fu, S.~Ermon, A.~Rudra, and C.~R{\'e}, ``FlashAttention: Fast and memory-efficient exact attention with IO-awareness,'' in \textit{Proc. NeurIPS}, 2022.

\bibitem{leviathan2023fast}
Y.~Leviathan, M.~Kalman, and Y.~Matias, ``Fast inference from transformers via speculative decoding,'' in \textit{Proc. ICML}, 2023.

\bibitem{chen2023accelerating}
C.~Chen, S.~Borgeaud, G.~Irving, J.-B. Lespiau, L.~Sifre, and J.~Jumper, ``Accelerating large language model decoding with speculative sampling,'' \textit{arXiv:2302.01318}, 2023.

\bibitem{dettmers2023qlora}
T.~Dettmers, A.~Pagnoni, A.~Holtzman, and L.~Zettlemoyer, ``QLoRA: Efficient finetuning of quantized language models,'' in \textit{Proc. NeurIPS}, 2023.

\bibitem{hu2022lora}
E.~J. Hu, Y.~Shen, P.~Wallis, Z.~Allen-Zhu, Y.~Li, S.~Wang, L.~Wang, and W.~Chen, ``LoRA: Low-rank adaptation of large language models,'' in \textit{Proc. ICLR}, 2022.

\bibitem{pope2023efficiently}
R.~Pope, S.~Douglas, A.~Chowdhery, J.~Devlin, J.~Bradbury, A.~Levskaya, J.~Heek, K.~Xiao, S.~Agrawal, and J.~Dean, ``Efficiently scaling transformer inference,'' in \textit{Proc. MLSys}, 2023.

\bibitem{kury2023clinical}
F.~Kury, A.~Butler, C.~Yuan, L.~Fu, and Y.~Sun, ``Automating clinical trial eligibility screening with natural language processing,'' \textit{JAMIA}, 2023.

\bibitem{vaswani2017attention}
A.~Vaswani, N.~Shazeer, N.~Parmar, J.~Uszkoreit, L.~Jones, A.~N. Gomez, {\L}.~Kaiser, and I.~Polosukhin, ``Attention is all you need,'' in \textit{Proc. NeurIPS}, 2017.

\bibitem{bai2024qwen}
J.~Bai et~al., ``Qwen2.5 technical report,'' \textit{arXiv:2412.15115}, 2024.

\end{thebibliography}

% ============================================================================
% Appendix
% ============================================================================

\appendix
\section{AI-Use Declaration}

This project utilized AI-assisted development tools during the coding and documentation process. Specifically:

\begin{itemize}
\item \textbf{Code development}: AI coding assistants were used to help structure the codebase, implement standard patterns (data loading, training loops), and debug issues. All algorithmic implementations (KV-cache management, tiled attention with online softmax, speculative decoding) were reviewed for correctness and understanding.

\item \textbf{Documentation}: AI tools assisted in drafting documentation, docstrings, and portions of this report. All content was reviewed and verified for accuracy.

\item \textbf{Nature of assistance}: The AI tools served as a programming aid analogous to an IDE with advanced autocomplete. The architectural decisions, algorithm designs, experimental methodology, and analysis were directed by the author.
\end{itemize}

All code, experiments, and analysis in this project represent the author's understanding and have been verified for correctness.

\end{document}
